% Above: the measurement. Below: the open-loop step through the whole path.
\begin{tikzpicture}[font=\sffamily\scriptsize, x=1cm, y=1cm]

% ================================================================ the sweep
\node[chlabel, anchor=west] at (0,3.6) {the dead-time sweep: asked for, encoded, and measured on the pins};

\foreach \i/\ask/\enc/\meas in {0/{100}/{100}/{103}, 1/{250}/{250}/{253},
                                2/{500}/{507}/{510}, 3/{1000}/{1014}/{1017},
                                4/{2000}/{2028}/{2031}}{
  \node[regcell, minimum width=17mm, text width=15mm, anchor=south west, fill=white]
    at (\i*1.85,2.5) {\ask\,ns\\{\tiny asked}};
  \node[regcell, minimum width=17mm, text width=15mm, anchor=south west, fill=kercol]
    at (\i*1.85,1.8) {\enc\,ns\\{\tiny encoded}};
  \node[regcell, minimum width=17mm, text width=15mm, anchor=south west, fill=hwcol]
    at (\i*1.85,1.1) {\meas\,ns\\{\tiny measured}};
}
\node[tlabel, anchor=west] at (9.4,2.8) {the request};
\node[tlabel, anchor=west] at (9.4,2.1) {what the piecewise encoding gives};
\node[tlabel, anchor=west] at (9.4,1.4) {what the two capture channels see};
\node[tlabel, anchor=west, text=red!60!black] at (9.4,0.8)
  {a constant +3 ns: the propagation difference, reported and not subtracted};

% ================================================================ the step
\node[chlabel, anchor=west] at (0,0.1) {the open-loop step: one duty command, two positions};

\draw[taxis] (0,-3.4) -- (9.0,-3.4);
\draw[taxis] (0,-3.4) -- (0,-0.6);
\node[font=\sffamily\tiny, anchor=east, rotate=90] at (-0.35,-2.0) {position};
\node[tlabel, anchor=north] at (4.5,-3.55) {two seconds after a duty step of 0.3};

% the modelled position: a smooth second-order rise
\draw[signal, line width=1pt]
  (0,-3.35) .. controls (2.4,-3.2) and (4.0,-1.6) .. (9.0,-1.0);
\node[tlabel, anchor=west] at (9.1,-1.0) {modelled position};

% the decoded position: the same, quantised and one period late
\draw[signalb, line width=0.9pt]
  (0.25,-3.35) .. controls (2.65,-3.2) and (4.25,-1.62) .. (9.0,-1.06);
\node[tlabel, anchor=west] at (9.1,-1.45) {decoded position, through the real path};

\draw[tspan] (0,-3.62) -- node[tlabel]{one control period of latency} (0.25,-3.62);

\node[umlnote, text width=52mm, anchor=north west] at (0,-4.1)
  {The two curves differ by one control period of latency and by the decoder's
   quantisation, and that difference is the point of routing the model through
   real hardware. A controller tuned against the upper curve alone would be
   tuned against a signal no encoder produces.};
\node[umlnote, text width=52mm, anchor=north west] at (5.6,-4.1)
  {Every number on the lower half of this figure is modelled and labelled so.
   The motion is fictional; the latency, the quantisation and the decode path
   are real, and they are what chapter 16's controller will actually meet.};
\end{tikzpicture}
